离线评估给出 AUC \(0.88\)，上线之后是 \(0.74\)。这样的落差在数据团队里常见到不值得惊讶，而它的成因几乎总落在同一处：某个本该在划分之后做的动作，提前发生在了全量数据上。这一单元处理的正是这类问题——手上只有一份数据，却要回答``换一批数据会怎样''，于是只能把想问的那种重复在这份数据上重放一遍。重放的方式有三种，各自替换掉了不同的东西，也各自会在不同的地方失灵。

第一种替换的是分布。抽样分布定义在那些没有发生的重复上，而重复需要从真实分布里再抽一次；既然分布未知，就用手上这份数据的经验分布顶上，有放回地重抽、重算，几千次之后得到一条可以量的曲线。整套方法的全部内容就是这一次顶替，因此它在经验分布还不像真实分布的地方——样本很小、尾巴很重、统计量盯着最大值这类边界量——会失手，而且失手时留有可以看见的症状。这是\hyperref[sec:11-Mathematical-Statistics/06-Resampling-and-Model-Selection/01]{``Bootstrap 近似估计量的抽样波动''}。

第二种替换的是标签。若处理毫无效果，每个单位的取值与它被分到哪一组无关，于是把标签打乱重算统计量，就得到了原假设下的参照分布，整个过程只有计数、没有分布假设。它并没有取消假设，只是把一条换成了另一条：打乱必须保持联合分布不变，也就是可交换性。同一个用户的多条记录、按公司发放的灰度、时间序列、配对设计，都会让逐行打乱变成非法操作，正确的做法是让置换的单位与随机化的单位对齐。这是\hyperref[sec:11-Mathematical-Statistics/06-Resampling-and-Model-Selection/02]{``置换检验依赖可交换性''}。

第三种替换的是训练流程本身。交叉验证把数据切成若干折，轮流留一折出来考试，它给出的数评价的是一整条流水线在某个数据规模上的表现，而不是某个具体的参数向量。这句话有直接的操作后果：标准化、缺失填补、特征筛选、目标编码、过采样、早停、超参搜索，凡是读过数据的步骤都必须搬进每一折内部重做一遍，否则验证折的信息会从上游倒流回来，估计系统性偏乐观。这一篇也给出几种半小时内能做完、确实能抓出事故的检查。这是\hyperref[sec:11-Mathematical-Statistics/06-Resampling-and-Model-Selection/03]{``交叉验证评估的是完整训练流程''}。

最后一篇把这些读数用起来。预测误差可以拆成不可约噪声、平方偏差与估计方差，后两者随复杂度反向变化，于是正则强度有一个中间的最优位置，交叉验证的任务就是找到它。麻烦在于验证曲线是真实风险加了一层噪声，在若干候选里取最小值这个动作本身就会把读数拉低；试的次数越多，验证集被拟合得越彻底。补救办法只有一类，就是让最后的评估落在没参与过选择的数据上。这是\hyperref[sec:11-Mathematical-Statistics/06-Resampling-and-Model-Selection/04]{``偏差---方差、正则化与模型选择''}。

有一件事贯穿四篇：三种方法的实现都长得像``循环加随机抽样''，逻辑却互不替代。Bootstrap 重抽观测，量的是估计量的波动；置换重排标签，给的是原假设下的参照分布；交叉验证重跑流程，估的是泛化误差。它们成立所依赖的独立性与可交换性条件各不相同，把一段代码改个抽样方式，得到的数就不能再承担原来的解释。
